% Completed-I input statements only; no FT/C.13 proof.

% Completed I statement: I:def:finite-coordinates
\begin{definition}[The finite local coordinates]
\hypertarget{tag-PI07}{}\label{I:def:finite-coordinates}
\StatementRecord{I:def:finite-coordinates}{proved}{PI07}{}
At the order-four point choose $t-1=s^4$ so that on the positive
$t>1$ arc toward the cusp, $\tau=iy$ with $y>1$, and
\begin{equation}\hypertarget{tag-01T5}{}\label{inv:eq:A-s-branch}
\tau=i+\frac{i}{\pi\sqrt{E_4(i)}}s^2+O(s^4),\qquad
\operatorname{Im}\mu_s(0)>0.
\end{equation}
Replacing $s$ by $i^js$ changes $\mu_s(0)$ by $i^{-j}$.
For any other local coordinate $z$ the derivative is transported
by $\mu_z=\mu_s s_z$. At the order-three point use $r^3=t$,
with $r>0$ on the arc $0<t<1$ along which $\tau$ travels from
$\rho$ to $i$ on the unit circle.
\end{definition}

% Completed I statement: inv:A2.3
\begin{proposition}[The elliptic logarithm at the finite points]
\hypertarget{tag-01RW}{}\label{inv:A2.3}
\StatementRecord{inv:A2.3}{proved}{01RW}{\SeriesRef{fam:src:B.1.1}, \SeriesRef{fam:src:B.1.3}; \citet{DukeImamoglu2008}}
Let $[\mu]$ be the generator identified in
Proposition~\ref{fam:src:B.1.1} with
one of the two sections
$P_*=(0,\sqrt2t^2(t-1))$ and $-P_*$ of
$y^2=x^3-3t^3(t-1)x+2t^4(t-1)^2$.
Along the branch in \eqref{inv:eq:A-s-branch}, put $x_0=1-1/t$.
Then
\begin{equation}\hypertarget{tag-01SY}{}\label{inv:eq:A-DI-branch}
\mu=\frac{1-\tau}{2}
 +\frac{i\sqrt6}{3\pi}
 \frac{x_0^{1/4}\,{}_3F_2(1/3,2/3,1;3/4,5/4;x_0)}
 {{}_2F_1(1/12,5/12;1;1-x_0)}.
\end{equation}
The difference $\mu-(2-\tau)/3$ has exact order two at the
order-three point in the coordinate $r^3=t$.
\end{proposition}

% Completed I statement: inv:A2.1
\begin{proposition}[The complete raw period determinant]\hypertarget{tag-01RU}{}\label{inv:A2.1}\label{wb:prop:log-ks-divisor}\label{wb:prop:period-derivative-determinant}
\StatementRecord{inv:A2.1}{proved}{01RU}{the definitions and constructions stated in the proof}

On the marked upper half-plane,
\begin{equation}\hypertarget{tag-01T3}{}\label{inv:eq:A-modular-K}
6\mu_z^2-\tau_z\beta_z
=\tau_z^2\frac{E_4(\tau)^2}{E_6(\tau)}.
\end{equation}
The right side extends holomorphically through the apparent poles.
Its exact orders are two at $z_1$ and zero at $z_2$; it has no other
zero away from their translates.  The logarithmic connecting map
$\kappa_0:f_*\mathcal E\to\mathcal L\otimes G$ has divisor
$\operatorname{divisor}(\det\kappa_0)=p_1+p_2$.
\end{proposition}

% Completed I statement: inv:A2.2
\begin{proposition}[The finite period coefficients]\hypertarget{tag-01RV}{}\label{inv:A2.2}
\StatementRecord{inv:A2.2}{proved}{01RV}{\SeriesRef{inv:eq:A-s-branch}, \SeriesRef{inv:eq:A-modular-K}, \SeriesRef{inv:A2.3}}

In the coordinate \eqref{inv:eq:A-s-branch},
\begin{equation}\hypertarget{tag-01T4}{}\label{inv:eq:A-mu-value}
\mu_s(0)=i\,C_4,
\qquad
C_4=\frac{\sqrt{2/3}}{\pi E_4(i)^{1/4}}
=\frac{4\sqrt\pi}{3^{3/4}\Gamma(1/4)^2}>0.
\end{equation}
For $r^3=t$, with $r>0$ on the arc $0<t<1$ and $\tau$ travelling
from $\rho$ to $i$ along the unit circle, set
\[
\lambda=\frac{3e^{5\pi i/6}}{2\pi E_6(\rho)^{1/3}}.
\]
Then $\tau=\rho+\lambda r+O(r^2)$ and
\[
K_r=-\frac{3i\lambda}{2\pi}r^2+O(r^3).
\]
The exact expressions have the rational enclosures
\[
\begin{aligned}
0.2366093140807852&<C_4<0.2366093140807853,\\
0.0801022110875382&<\operatorname{Re}K_2<0.0801022110875383,\\
0.1387410994022233&<\operatorname{Im}K_2<0.1387410994022234.
\end{aligned}
\]
\end{proposition}

% Completed I statement: inv:C3.2
\begin{proposition}[The factor $12$ with a fixed collar orientation]\hypertarget{tag-01SF}{}\label{inv:C3.2}\label{wb:lem:bcpr-degree-one-input}
\StatementRecord{inv:C3.2}{proved}{01SF}{\SeriesRef{inv:eq:C-collar-sum}, \SeriesRef{inv:C3.1}}

Use the marked cohomology class $\gamma$ dual to
$\hat\gamma$, the analytic peripheral generators and the lift relation of
\cite[\S7.5, pp.~54--57, especially (7.10), p.~55]{AC2026}, and put
$p=12l_0-4l_1-3l_2=\pm1$.  Let $g_\gamma$ be the unique coherent
Leray class whose fibre image is $\gamma^{0,1}$.  Then
\begin{equation}\hypertarget{tag-01TE}{}\label{inv:eq:C-factor12}
\exp(12g_\gamma)=f^*\OO_{\Pone}(-p).
\end{equation}
For the standard data $(l_0,l_1,l_2)=(0,1,-1)$, this is
$\exp(12g_\gamma)=f^*\OO(1)$.  For the opposite value $p=1$,
the exponent representing $f^*\OO(1)$ is $-12g_\gamma$.
\end{proposition}
